\chapter{Conclusioni e limiti}
\section{Risultati principali}
Il lavoro ha costruito una catena coerente dalla meccanica non lineare alla valutazione robusta. Il manipolatore SEA è stato descritto mantenendo separate coordinate dei link e dei rotori, quindi linearizzato attorno all'equilibrio verticale superiore. La dinamica degli attuatori e il ritardo hanno portato la pianta nominale a 14 stati. L'insieme M4 combina variazioni parametriche reali e dinamica moltiplicativa non modellata, mentre M5 traduce i requisiti temporali in pesi frequenziali e scale fisiche.

La progettazione nominale mostra che M6 LQG, senza integratore, è il riferimento più efficace per il tracking: raggiunge il target di settling e mantiene overshoot e cross-coupling entro i limiti adottati. L'offset da disturbo costante rimane però non nullo. M7 dimostra che l'introduzione dell'internal model elimina tale offset, ma nella LQRI con Kalman M6 e integratore misurato studiata non consente di raggiungere il target di rapidità. La LQRI è implementata e stabile sul nominale, ma resta non conforme alle specifiche M5; non viene promossa a sostituto del riferimento M6.

Le sintesi $H_\infty$ M9 e M10 verificano l'equivalenza fra \texttt{mixsyn} e una generalized plant costruita esplicitamente. Entrambe restituiscono una norma pesata circa \HinfsynGamma, superiore a uno. M11 mostra come l'aumento dell'ordine delle strutture MIMO da A3 ad A5 riduca il divario dalla soluzione full-order; nessuna architettura soddisfa tuttavia l'obiettivo pesato. A5 viene usato per la successiva analisi robusta, non confuso con $K_{\infty}$.

La D--K iteration M13 produce $K_\mu$ di ordine \KmuOrder. La post-analisi M14, separata dalla sintesi, costituisce il risultato formale centrale: Robust Stability è \status{\KmuRSStatus} con margine [\KmuRSLower,\KmuRSUpper] sull'insieme lineare incerto M4. Nominal Performance e Robust Performance restano \status{FAIL} e \status{FAIL\_CERTIFIED}. L'intervallo RS è sotto uno: il risultato è negativo anche per la stabilità robusta e non eredita il certificato della baseline precedente.

Le campagne non lineari e Monte Carlo completano la lettura. M6 conserva il miglior tracking nominale, ma M16 osserva una perdita di stabilità nei Sample ID \MControllerUnstableSampleIds. $K_{\infty}$ e $K_\mu$ sono stabili rispettivamente in \KhinfStableSamples/\MonteCarloSamples\ e \KmuStableSamples/\MonteCarloSamples\ campioni, pur fallendo il tracking M5. La frazione empirica di $K_\mu$ non sostituisce M14 e non ne ribalta il risultato negativo. Sul modello non lineare nominale tutte le architetture mostrano regioni empiriche di convergenza significative; le stime search-box limited non sono trasformate in frontiere certificate.

\section{Decisione ingegneristica}
La verifica aggiuntiva di Lyapunov riguarda l'intera dinamica nominale aumentata
di ciascun controllore, non il solo sottosistema meccanico. Una maggiorazione
analitica del resto cubico fornisce tre regioni interne sufficienti, piccole
e conservative. Le ellissi numeriche pi\`u ampie e le mappe di traiettorie
non sono certificate. Questo risultato non estende M14 al non lineare incerto
e non sostituisce i requisiti prestazionali falliti.

Non emerge un controllore che domini tutte le metriche. La selezione finale è pertanto per ruolo:
\begin{itemize}
\item il ruolo di controllore robusto finale rimane non assegnato, per assenza di RS certificata;
\item M6 è il riferimento di prestazione nominale, per il tracking conforme;
\item $K_{\infty}$ è il comparatore metodologico full-order della famiglia $H_\infty$.
\end{itemize}
Questa assegnazione evita due semplificazioni scorrette: attribuire un certificato a $K_\mu$ perché proviene dalla sintesi robusta, oppure dichiarare M6 robusto perché supera quasi tutti i campioni. La conclusione documenta quale evidenza sostiene ogni ruolo e quale evidenza manca.

\section{Limiti noti}
% Tabella generata da build_m18_report_tables.m.
\begin{itemize}
\item Nessun controllore principale soddisfa tutte le specifiche M5
\item Kmu possiede RS formale, ma fallisce NP e RP in M14
\item M6 e Khinf non dispongono di una certificazione RS formale su M4
\item M16 fornisce evidenza empirica lineare e non una certificazione di robustezza
\item Le RAS M8/M15 sono empiriche; i limiti del box non sono frontiere dinamiche
\item uScale=12 N m è una normalizzazione; non è dichiarata alcuna saturazione
\item Il limite interno di Lyapunov vale sul nominale aumentato; il livello piu ampio campionato non e certificato
\item La maggiorazione e conservativa e valutata in double, senza aritmetica intervallare
\item Nessuna riduzione o nuova sintesi robusta nella revisione; sono aggiunte analisi nominali di Lyapunov e sensibilita Kalman
\end{itemize}


I limiti hanno conseguenze dirette. Senza dati di saturazione e rate limit, i picchi di comando non bastano a dimostrare l'implementabilità sull'hardware. Senza identificazione sperimentale, i range M4 restano scelte ingegneristiche motivate ma non intervalli statistici. Senza un modello non lineare incerto, la certificazione lineare e la simulazione non lineare rimangono complementari, non unificate. La Monte Carlo usa soltanto 40 campioni e non autorizza inferenze probabilistiche.

L'ordine \KmuOrder\ di $K_\mu$ aumenta memoria, costo computazionale e sensibilità numerica di un'eventuale implementazione discreta. Nessuna riduzione è stata applicata nella fase conclusiva, perché preservare approssimativamente la risposta nominale richiederebbe una nuova verifica completa dei margini. Infine, le regioni empiriche di attrazione dipendono da box, griglia, orizzonte e criterio di convergenza; per M6 e $K_\mu$ la vera frontiera non è stata localizzata.

\section{Sviluppi futuri}
Per ampliare le regioni interne di Lyapunov si potrebbero studiare coordinate
meglio scalate e maggiorazioni meno conservative, senza interpretare il
livello numerico campionato come prova. Una verifica con aritmetica
intervallare sarebbe necessaria per certificare anche gli arrotondamenti
del calcolo; non \`e stata eseguita in questa revisione.

Un seguito sperimentale dovrebbe partire dalla caratterizzazione degli attuatori: limiti di coppia e velocità, saturazioni, rate limit, attriti e ritardi misurati. Questi dati consentirebbero di sostituire le sole scale prestazionali con vincoli verificabili e di aggiornare l'incertezza dinamica. Una campagna di identificazione su più configurazioni potrebbe inoltre valutare se le variazioni di rigidezza e inerzia adottate in M4 descrivano correttamente l'impianto.

Sul piano del controllo, lo studio mirato a tre pesi del Capitolo~8
costituisce un primo confronto motivato ma non risolve i target nominali:
i suoi nuovi $W_1,W_2$ sono stati adottati nella catena, che è stata rivalutata senza riusare i precedenti certificati. Un'eventuale configurazione
futura da adottare richiederebbe una nuova iterazione completa, senza
cancellare gli esiti presenti. Architetture 2-DoF potrebbero separare tracking e reiezione dei disturbi; l'anti-windup diventerebbe necessario una volta definiti i limiti di attuazione. Una riduzione di $K_\mu$ dovrebbe essere accompagnata da una post-analisi M14 completa, dalle simulazioni non lineari e da una nuova Monte Carlo sugli stessi criteri.

Infine, l'estensione più rilevante sarebbe una validazione non lineare incerta, possibilmente abbinata a prove hardware-in-the-loop. Tale fase collegherebbe la robustezza formale locale alle dinamiche fisiche oltre l'equilibrio. Fino ad allora, il risultato raggiunto resta deliberatamente preciso: la configurazione corrente dispone di un riferimento nominale con tracking conforme, ma non di un controllore principale con RS certificata; documenta esplicitamente questa limitazione e le prestazioni non soddisfatte.
